% Part: sets-functions-relations
% Chapter: functions
% Section: kinds

\documentclass[../../../include/open-logic-section]{subfiles}

\begin{document}

\olfileid[ar]{sfr}{fun}{kin}
\olsection{أنواع الدوال}

\begin{explain}
ولكثرة ورود بعض أصناف الدوال ننتفع بتصنيف يميزها.

فلنبدأ بما تبلغ قيمه كل عنصر من المجال المقابل. وتسمى هذه الدوال !!{surjective}، وصورتها في \olref{fig:surjective}.

\begin{figure}
  \olasset{assets/diagrams/surjective.tikz}
  \caption{الدالة !!^a{surjective} تبلغ بقيمها كل !!{element} من المجال المقابل.}
  \ollabel{fig:surjective}
\end{figure}
\end{explain}

\begin{defn}[دالة !!^{surjective}]
الدالة ${\OLId{latin-ordinary}{f}} \colon {\OLId{latin-looped}{A}} \rightarrow {\OLId{latin-looped}{B}}$ \emph{!!{surjective}} إذا وفقط إذا
اتحد ${\OLId{latin-looped}{B}}$ ومدى~${\OLId{latin-ordinary}{f}}$؛ أي لكل ${\OLId{latin-ordinary}{y}} \in {\OLId{latin-looped}{B}}$ يوجد ${\OLId{latin-ordinary}{x}} \in {\OLId{latin-looped}{A}}$ واحد على الأقل
يحقق~${\OLId{latin-ordinary}{f}}({\OLId{latin-ordinary}{x}}) = {\OLId{latin-ordinary}{y}}$. وصورته الرمزية:
\[
  (\forall {\OLId{latin-ordinary}{y}} \in {\OLId{latin-looped}{B}})(\exists {\OLId{latin-ordinary}{x}} \in {\OLId{latin-looped}{A}}){\OLId{latin-ordinary}{f}}({\OLId{latin-ordinary}{x}}) = {\OLId{latin-ordinary}{y}}.
\]
وتسمى مثل هذه الدالة !!a{surjection} من ${\OLId{latin-looped}{A}}$ إلى ${\OLId{latin-looped}{B}}$.
\end{defn}

\begin{explain}
فإثبات أن ${\OLId{latin-ordinary}{f}}$ هي !!a{surjection} يكون بإظهار أن كل شيء في المجال
المقابل لـ${\OLId{latin-ordinary}{f}}$ يحصل قيمةً ${\OLId{latin-ordinary}{f}}({\OLId{latin-ordinary}{x}})$ لمدخل ${\OLId{latin-ordinary}{x}}$ ما.

وكل دالة \emph{تستحث} !!a{surjection}. خذ ${\OLId{latin-ordinary}{f}} \colon {\OLId{latin-looped}{A}} \to {\OLId{latin-looped}{B}}$،
واجعل ${\OLId{latin-ordinary}{f}}' \colon {\OLId{latin-looped}{A}} \to \ran{{\OLId{latin-ordinary}{f}}}$ معرفة بـ${\OLId{latin-ordinary}{f}}'({\OLId{latin-ordinary}{x}}) = {\OLId{latin-ordinary}{f}}({\OLId{latin-ordinary}{x}})$.
فـ$\ran{{\OLId{latin-ordinary}{f}}}$ \emph{بتعريفه} $\Setabs{{\OLId{latin-ordinary}{f}}({\OLId{latin-ordinary}{x}}) \in {\OLId{latin-looped}{B}}}{{\OLId{latin-ordinary}{x}} \in {\OLId{latin-looped}{A}}}$،
ومن ثم ${\OLId{latin-ordinary}{f}}'$ هي !!a{surjection}.
\end{explain}

\begin{explain}
ولكل مدخل مخرج واحد في أي دالة. وقد تزيد الدالة على ذلك بأن
لا تجمع مدخلين مختلفين على مخرج واحد؛ فتسمى !!{injective}،
وتمثيلها في \olref{fig:injective}.
\begin{figure}
  \olasset{assets/diagrams/injective.tikz}
  \caption{الدالة !!^a{injective} تمنع اتفاق القيمة لوسيطتين مختلفتين.}
  \ollabel{fig:injective}
\end{figure}
\end{explain}

\begin{defn}[دالة !!^{injective}]
الدالة ${\OLId{latin-ordinary}{f}} \colon {\OLId{latin-looped}{A}} \rightarrow {\OLId{latin-looped}{B}}$ \emph{!!{injective}} إذا وفقط إذا
كان لكل ${\OLId{latin-ordinary}{y}} \in {\OLId{latin-looped}{B}}$ عنصر ${\OLId{latin-ordinary}{x}} \in {\OLId{latin-looped}{A}}$ واحد على الأكثر يحقق~${\OLId{latin-ordinary}{f}}({\OLId{latin-ordinary}{x}}) = {\OLId{latin-ordinary}{y}}$.
وتسمى حينئذ !!a{injection} من ${\OLId{latin-looped}{A}}$ إلى~${\OLId{latin-looped}{B}}$.
\end{defn}

\begin{explain}
ولإثبات أن ${\OLId{latin-ordinary}{f}}$ هي !!a{injection} خذ !!{element}s ${\OLId{latin-ordinary}{x}}$ و${\OLId{latin-ordinary}{y}}$ من مجال
${\OLId{latin-ordinary}{f}}$، وبيّن أن ${\OLId{latin-ordinary}{f}}({\OLId{latin-ordinary}{x}})={\OLId{latin-ordinary}{f}}({\OLId{latin-ordinary}{y}})$ لا يكون إلا مع ${\OLId{latin-ordinary}{x}}={\OLId{latin-ordinary}{y}}$.
\end{explain}

\begin{ex}
خذ الثابتة ${\OLId{latin-ordinary}{f}}\colon \Nat \to \Nat$ المحددة بـ${\OLId{latin-ordinary}{f}}({\OLId{latin-ordinary}{x}}) = {\OLMathNumeral{1}}$؛ فهذه
لا هي !!{injective} ولا هي !!{surjective}.

وأما الهوية ${\OLId{latin-ordinary}{f}}\colon \Nat \to \Nat$، حيث ${\OLId{latin-ordinary}{f}}({\OLId{latin-ordinary}{x}}) = {\OLId{latin-ordinary}{x}}$،
فهي !!{injective} و!!{surjective} معًا.

وأما اللاحق ${\OLId{latin-ordinary}{f}} \colon \Nat \to \Nat$، حيث ${\OLId{latin-ordinary}{f}}({\OLId{latin-ordinary}{x}}) = {\OLId{latin-ordinary}{x}}+{\OLMathNumeral{1}}$،
فهي !!{injective} دون !!{surjective}.

وأما الدالة ${\OLId{latin-ordinary}{f}} \colon \Nat \to \Nat$ ذات التعريف:
\[
  {\OLId{latin-ordinary}{f}}({\OLId{latin-ordinary}{x}}) =
  \begin{cases}
    \frac{{\OLId{latin-ordinary}{x}}}{{\OLMathNumeral{2}}} & \text{إذا كان ${\OLId{latin-ordinary}{x}}$ زوجيًا} \\
    \frac{{\OLId{latin-ordinary}{x}}+{\OLMathNumeral{1}}}{{\OLMathNumeral{2}}} & \text{إذا كان ${\OLId{latin-ordinary}{x}}$ فرديًا.}
  \end{cases}
\]
فهي !!{surjective} غير !!{injective}.
\end{ex}

\begin{explain}
ويكثر احتياجنا إلى الدوال التي تكون !!{injective} و!!{surjective} معًا؛
واسمها !!{bijective}، وصورتها في \olref{fig:bijective}.
وتسمى !!^{bijection}s أيضًا \emph{مراسلات واحد إلى واحد}، إذ يقابل
كل عنصر من المجال المقابل فيها عنصر واحد من المجال لا غير.
\begin{figure}
  \olasset{assets/diagrams/bijective.tikz}
  \caption{الدالة !!^a{bijective} تجعل لكل عنصر من المجال المقابل مقابِلًا واحدًا في المجال.}
  \ollabel{fig:bijective}
\end{figure}
\end{explain}

\begin{defn}[!!^{bijection}]
الدالة ${\OLId{latin-ordinary}{f}} \colon {\OLId{latin-looped}{A}} \to {\OLId{latin-looped}{B}}$ \emph{!!{bijective}} إذا وفقط إذا كانت !!{surjective} و!!{injective} معًا. وحينئذ تسمى !!a{bijection}
من ${\OLId{latin-looped}{A}}$ إلى~${\OLId{latin-looped}{B}}$، أو بين ${\OLId{latin-looped}{A}}$ و~${\OLId{latin-looped}{B}}$.
\end{defn}

\end{document}
